\section{The cubic coefficient}\label{sec:cubic}

\begin{theorem}\label{thm:cubicdisk}
For every Young diagram the two vacuum-normalized disk coefficients of
exchanged degree three satisfy $T_3=U_3$ in $\mathbb Q(p,c,x)$.
\end{theorem}

Write $M$ for the chronological moment, and $G$ for the moment obtained by
putting the negatives before the positives, preserving the order within
each sign. The argument begins by removing the scalar part of a change of order. We then prove the adjacent-slice permutation identity and compute the residues produced on restoring the original order.

\subsection{Removing the longer gaps}
Put $\mathcal C_h=l_1 C_h$ and $a_0=(1-s)^2(1-1/p)/(1-s/p)$, so
$l_1\gamma a_0=-1$. The scalar term in the difference of orders of a
positive at $k$ and a later negative at $j$ is
\[
 b(k,j)=a_0(1-s^{j-k})/(1-s)\quad (k<j),\qquad b(k,j)=0\quad(k\ge j).
\]
In a word of ordinary fields subtract this scalar at every inverted
pair by inclusion and exclusion. More precisely, set
\[
 M^\circ(w)=\sum_D(-1)^{|D|}\prod_{(k,j)\in D}b(k,j)\,M(w\setminus D),
\]
where $D$ ranges over matchings and $w\setminus D$ deletes both ends
of every matched pair. A single deletion in a signed sum joins
two intervals, since
\[
 b(i,j)-s b(i,j-1)-b(i+1,j)+s b(i+1,j-1)=a_0{\bf1}_{j=i+1}.
\]
The two untouched ends give the current of the joined interval. The
relative scalar is $l_1\gamma a_0=-1$. Consequently the six directed
single edges on three labeled insertions contribute $M(\mathcal C_1\mathcal C_2)$ to
$M^\circ(\mathcal C_1^3)/6$. The six directed paths of length two contribute
$M(\mathcal C_3)$. This proves that the whole third coefficient is
\begin{equation}\label{eq:cubic-moment}
 \tfrac16 M^\circ(\mathcal C_1^3)+\tfrac12 l_2 M(\mathcal C_1\psi^2 C_1)
                         +\tfrac13 l_3 M(\psi^3 C_1).
\end{equation}
Here $\psi^j$ is the $j$th Adams operation on the diagonal eigenvalues.
No subtraction is made in the second summand. The two fields of the
last summand are already in grouped order.

\subsection{Permuting the slice parameters}
We give explicitly the identity used to sum the slices. The two
alphabets $A,B$ below label positive and negative residues respectively.
The label of a positive variable of value $v$ is $v/s$, and that of
a negative variable of value $w$ is $xw$. Passing a slice parameter $X$
in the forward positive or backward negative direction multiplies
the corresponding one-body weight by
$r_z(X)=(X-z)/(X-sz)$. Initially take distinct indeterminates in each
alphabet. Put $r_J(X)=\prod_{z\in J}r_z(X)$.

For one sign let $(q,V)=(p,s)$ or $(1/p,1/s)$ as appropriate and write
\[
 K(t)=\frac{(1-t)(1-qt/V)}{(1-qt)(1-t/V)},\qquad
 E(J)=\prod_{\{z,w\}\subset J}
       \frac{-V(z-w)^2}{(Vz-w)(Vw-z)} .
\]
We also write $K_{q,V}$ to display the parameters, and
$K(I,J)=\prod_{z\in I,w\in J}K(w/z)$. Let $S(J)$ be the sum, over
the orders on $J$, of the products of the corresponding ordered pair
factors. For a power sum put
\[
 P(J)=
 \begin{cases}
 1,&|J|=1,\\
 S(J)-2E(J),&|J|=2,\\
 S(J)-3\sum_{u\in J}K(\{u\},J\setminus\{u\})E(J\setminus\{u\})+3E(J),
   &|J|=3 .
 \end{cases}
\]
These are Newton's formulas. The Cauchy determinant of an elementary
operator, divided by the singleton factors, is precisely $E(J)$
\cite[Props.~2.6--2.7]{Ko}.

For a partition $\nu=(j_1,\ldots,j_m)$ of $r$ in a fixed order, define
$T^\pm_\nu(A;X)$ as follows, using the appropriate sign throughout.
Sum over ordered decompositions $A=J_1\sqcup\cdots\sqcup J_m$ with
$|J_i|=j_i$, and over the two cuts for each group. Start with
$\prod_iP(J_i)$. The displaced cut of a positive group contributes
$-r_{J_i}(X)$, that of a negative group $-s^{j_i}r_{J_i}(X)$.
Multiply the pair factors $K(J_i,J_j)$ in cut order. Thus at one
center the positive at the center precedes a displaced positive, and
a displaced negative precedes the negative at the center. Use list
order at ties. For $|A|=|B|=r$ put
\begin{equation}\label{eq:cubic-density}
 D_X(A,B)=\sum_{\nu\vdash r} w_\nu T^+_\nu(A;X)T^-_\nu(B;X),
 \qquad w_\nu=\prod_{j\ge1}\frac{1}{m_j!}
              \left(\frac{l_j}{j l_1^j}\right)^{m_j},
\end{equation}
where $m_j$ is the multiplicity of $j$; the empty entry is one.
This is the contribution of a single slice after removal of its
one-body factors, the common opposite-sign contractions and
$(l_1\gamma)^r$. Indeed the scaled power sum has relative coefficient
$(V^j-1)/(V-1)^j$. The product for the two signs, times
$-s^j/(s^j-1)^2$, is $\gamma^j$.

For functions of paired, equally large alphabets let $*$ denote the
product obtained by summing over ordered decompositions
$(A,B)=(A_1,B_1)\sqcup(A_2,B_2)$ with $|A_i|=|B_i|$, inserting
$K_+(A_1,A_2)K_-(B_1,B_2)$ between the factors. The signs specify the
two choices of $(q,V)$. Write $R_X^+,R_X^-$ for multiplication by
$r_A(X),r_B(X)$, respectively. Pair factors make $*$ associative,
and $R_X^\pm$ preserve it.
\begin{lemma}\label{lem:cubic-slices}
In degrees at most three,
\begin{equation}\label{eq:cubic-commute}
 (R_Y^-D_X)*(R_X^+D_Y)=(R_X^-D_Y)*(R_Y^+D_X).
\end{equation}
Consequently the sum of the slice assignments is symmetric in their
parameters, after extracting the one-body weights at the two outer cuts.
\end{lemma}

\begin{proof}
The factors across the two slices are just the transports past $X$
for positives on the second slice and past $Y$ for negatives on the
first. Earlier centers precede later centers in the same-sign order,
including ties at the displaced cuts. The transports from outside the
pair depend only on its combined alphabets. Hence the second assertion
follows by adjacent transpositions.

Here are finite interpolation details for the new homogeneous part
of degree three. For each sign, index ten entries by pairs of partitions
$(\nu,\mu)$ with $|\nu|+|\mu|=3$. The positive entry $F^+_{XY}$ is
\[
 \sum_{I\sqcup J=A,\ |I|=|\nu|}
       r_J(X)K_+(I,J)T^+_\nu(I;X)T^+_\mu(J;Y).
\]
The negative entry $F^-_{XY}$ uses $B$, $K_-$, $T^-$ and $r_I(Y)$
in place of $r_J(X)$. Put $w_{\nu\mu}=w_\nu w_\mu$. The two sides of
\eqref{eq:cubic-commute} agree exactly when the rows
\[
 (F^+_{XY},-F^+_{YX}),\qquad
 (w_{\nu\mu}F^-_{XY},w_{\nu\mu}F^-_{YX})
\]
are orthogonal. The index $(\nu,\mu)$ runs through each of the two
ten-entry blocks. The following finite rational calculation suffices.
Take $Y=-1$, $X=-2,\ldots,-8$, $2\le p\le19$ and $21\le s\le50$,
with integral $X,p,s$. For each sign use every triple from the thirteen
primes from $11$ to $59$ as its labels. Form the two $286$ by $20$
rational matrices with the displayed rows. At each of the $3780$
choices of $(X,p,s)$ a row basis of the first matrix times the
transpose of the second is zero. The supplementary computation
implements these sums over subsets and uses rational row reduction.

For each unordered same-sign pair $z,w$ clear
$(z-pw)(w-pz)(z-sw)(w-sz)$. Each pair occurs in at most one ordered
or elementary factor of each summand. The resulting pair contribution
has degree at most two in each of $p,s$, four in each label. Clear also
$(X-sz)(Y-sz)$ for each label. The degrees are thus at most ten per
label and six per slice parameter. Negative displaced cuts supply at
most $s^3$. The only extra scalars are
\[
 \frac{l_2}{l_1^2}=\frac{(s+p)(1-p)}{(s-p)(1+p)},\qquad
 \frac{l_3}{l_1^3}=\frac{(s^2+sp+p^2)(1-p)^2}
                           {(s-p)^2(1+p+p^2)} .
\]
At most one occurs per summand. Clearing
$(s-p)^2(1+p)(1+p+p^2)$ gives final bounds $17$ in $p$ and $29$
in $s$. All grid denominators are nonzero. Symmetry in each alphabet
gives every ordering of its three distinct labels. With the other
two fixed, eleven nodes remain for each label; successively interpolating
therefore removes the distinct-node restriction and proves the identity
in the six labels. The seven slice values suffice for $X$.
Homogeneity under common rescaling restores $Y$, and the $18$ by
$30$ character grid finishes the proof. The lower degrees follow from Lemma~\ref{lem:twoslicetwo} and its degree-one calculation.
\end{proof}

\subsection{The change of order}
At a step $h$ write $\tau_h=Z_{h+1}/x$, which is $Z_h$ on an H step
and $Z_h/c$ on a D step. In the interchange of a positive at cut $k$
with a negative at $j>k$ a boundary residue is indexed by $k\le h<j$.
The two fields then have values $z_h=s Z_h$, $w_h=\tau_h$. Put
$S(v)=R(v)/R(cv)$ and
\[
 \begin{split}
 W_h(k,j)&=\prod_{k\le t<j,\ t\ne h} R(Z_h/Z_t),\\
 \Phi_{h,j}&=
 \begin{cases}
 \prod_{d\ge j,\ {\rm D}} S(Z_h/Z_d),&h\ {\rm H},\\
 \prod_{b<j,\ {\rm H}} S(Z_h/(cZ_b)),&h\ {\rm D}.
 \end{cases}
 \end{split}
\]
For chronological minus grouped order the residue has weight
$-a_0W_h(k,j)\Phi_{h,j}$. To check its action on the other fields,
expand the mixed kernel in the two orders. Its difference of expansions
is supported at $z=sw,pw$. At $sw$ the creation exponential of the
pair cancels and only contractions to its right remain; at $pw$ the
annihilation exponential cancels and only contractions to its left
remain. For a remaining field of sign $e=\pm1$ and value $u$, these
products are, respectively,
\[
 U_e(t)=K_{p,s^e}(t/s)K_{1/p,s^e}(t),\quad t=u/Z_h,\qquad
 V_e(t)=K_{p^e,s}(pt)K_{p^e,1/s}(t),\quad t=Z_h/(cu).
\]
The mixed kernel is unchanged as a rational function upon reversing
both the signs' order and the spectral ratio. Thus these same
products may be obtained by retaining all cross contractions in the
given same-sign order, even on the side where their product is one.
The two endpoints of the fused contour give $b(k,j)$ times the
moment with the pair deleted. The contour calculation may first be made
with independent boundary parameters, the horizontal ones sufficiently
large and the downward ones sufficiently small relative to all the
nested field circles. Its rational residues can then be specialized
to the $Z_h$ of the diagram.

\begin{lemma}\label{lem:cubic-poles}
After the scalar deletions in \eqref{eq:cubic-moment}, the strictly
inverted fields have either one or two boundary residues as above.
The remaining nonconstant choices are ordinary boundary residues and
their successive $p$-shifts, with one possible additional choice. At
an H residue it is a second negative with value $Z_h/c$; at a D
residue it is a second positive with value $sZ_h/c$. Either additional
choice occurs with just one inverted pair.
\end{lemma}

\begin{proof}
Consider first the ordinary cube. Number the fields of each sign in
their actual cut order, taking negatives first at ties. Every center
has its negative before its own positive, so each prefix contains at
least as many negatives as positives. The only candidate inversions are
$a=P_1N_2$, $d=P_2N_3$ and $b=P_1N_3$. Those among $a,d$ present
are adjacent disjoint pairs. Expand each as reversed order plus its
difference; after the two reversals only $b$ can require a swap.
There is one single difference for each inverted pair, and a double
difference precisely when both $a$ and $d$ occur. In the three possible
single differences the remaining words have orders
\[
 N_1[a]N_3P_2P_3,\qquad N_1N_2[b]P_2P_3,\qquad
 N_1N_2P_1[d]P_3 .
\]
In a double difference the intervals are disjoint since $N_2$
precedes $P_2$. The remaining fields are $N_1,P_3$.
In inclusion and exclusion, a matching using exactly one end of the
interchanged pair gives the same deleted word in both orders. The
matching of the pair itself subtracts its endpoint constant. Applying
this twice subtracts the two single constants and adds back their
product in the double difference. The resulting two boundary residues
retain their mutual contraction. For an earlier H and a later D it is
$U_+(pt)U_-(t)$, $t=Z_g/(cZ_h)$, where $g$ is the D step; otherwise
it is one. In the first case the outer H contour and inner D contour
move apart, so their mutual poles are never crossed.

In a single difference shrink the remaining positive contours from last
to first and enlarge the negatives from first to last. Their mutual
mixed poles stay on the other side. Repeating an ordinary boundary
selection gives its next $p$-shift by
\[
 R(t)K_{1/p,1/s}(t)=R(t/p)=R(t)K_{p,s}(1/t).
\]
At most one negative remains to the right of an H residue. At its
cut the one-body factor at $h$ and the two contractions with the pair
give, for $t=u/Z_h$,
\[
 R(t)U_-(t)=
 \frac{(1-t/s)(1-pt)(1-t/c)}
      {(1-ct)(1-t/p)(1-t)} .
\]
Thus the three choices at this scale have $t=1,p,1/c$. The other
negative in this word has cut at most $k$ and cannot also select $h$.
At a D residue at most one positive remains on the left. With
$t=u/(sZ_h)$ the product is
\[
 R(t)K_{p,s}(1/t)K_{p,1/s}(1/(pt)).
\]
It has the same three possible poles. The opposite orientations lie
outside the selected direction of contour motion. At distinct steps
the spectral monomials have distinct $x$-degrees.

The coincident choices $t=1$ cancel in pairs. On H the dressing of the
later negative at coincidence is $-1$, while the earlier negative is
undressed when the later one is fused. For the two eligible cuts $j_1,j_2$
the $j$-dependent part of the fused weight and the unfused H residue at
$Z_h$ have the same backward ratio $r_{x Z_h}(Z_j)$. Their products
agree upon exchanging $j_1,j_2$. All other factors have the same
evaluated variables and the same order for each sign. On D the
corresponding common forward ratio is $r_{Z_h}(Z_k)$ and the earlier
positive's dressing is again $-1$. This also gives cancellation at a
tie between two equally eligible fields. The endpoint weights of an
unfused ordinary positive and negative are $1$ and $s^j$ respectively.
Each contraction with the deleted endpoint tends to one. The primary
inequalities and the order of all the remaining fields are unchanged
when this field is deleted. In the sum over the candidate inversions
each such labeled inversion occurs once, so the two cuts of the
endpoint contribute $1-1$ or $s^i-s\,s^{i-1}$. For two inversions
the unfused fields are exterior and the same cancellation applies.

For the mixed summand apply $p_2=e_1^2-2e_2$ to the member of
the inverted pair coming from the second Adams operation.
Write $E_V(t)=-V(1-t)^2/((V-t)(Vt-1))$ for the two-variable
determinant factor, and put
\[
 S_{q,V}(t)=K_{q,V}(t)+K_{q,V}(1/t)-2E_V(t),\qquad
 \rho_V=(V^2-1)/(V-1)^2 .
\]
When the higher field is negative, each primary boundary residue is
multiplied by
\begin{equation}\label{eq:cubic-mixed}
 \rho_{1/s}\,\CT_u N_j(u)
 K_{p,1/s}(u/z_h)S_{1/p,1/s}(u/w_h).
\end{equation}
For a higher positive the multiplier is
$\rho_s\CT_u P_k(u)K_{1/p,s}(u/w_h)S_{p,s}(u/z_h)$.
The exterior negative and positive at the two centers may first be
evaluated on their prefix and suffix. Each chosen exterior root is
outside $[k,j)$, and its finite shift replaces that root by its
shifted value in the one-body factors. Equivalently, one retains
its regular contractions with the middle variables. The term with
no shift cancels in the corresponding signed difference.

Here is the contour order in \eqref{eq:cubic-mixed}. For two elementary
negatives take $0<|p|<|s|<1$ and equal radii for the two negative
variables. The determinant denominators $w_i-sw_j$ have their
prescribed expansions on these circles. Move the ordinary probe past
both. The two choices of the variable it meets cancel $1/2!$. At
$z=sw,pw$, with $t=u/w$ for the remaining variable, the respective
products $E_{1/s}(t)K_{p,1/s}(t/s)$ and $E_{1/s}(t)K_{p,1/s}(t/p)$
are
\[
 \frac{s(t-1)(pt-1)}{(st-1)(pt-s)},\qquad
 \frac{s(t-p)(t-1)}{(t-s)(st-p)} .
\]
The first has all its poles outside the unit circle, so $w$ can move
outwards, and the second has all its poles inside, so $w$ can move
inwards. Each gives one at the corresponding endpoint. For elementary
positives take $1<|s|<|p|$ and replace $t$ by $1/t$, moving the negative
probe. Both calculations continue rationally. For $e_1^2$ cross the
two ordinary fields successively. The two possibilities give
$K(t)+K(1/t)$ with the same mixed factor. Their two endpoint constants
cancel those of $2e_2$. This proves the stated contour formulas.

The function $S_{q,V}$ has a double zero at $1$ and is zero at $0$ and
$\infty$. At H set $t=u/Z_h$ in \eqref{eq:cubic-mixed}. The same-sign
denominators allow $p,1/p,s,1/s$. The mixed numerator cancels $1/p,s$,
and the one-body zero cancels $1/s$. The double zero at $1$ cancels
both the coincident poles. Besides ordinary H steps, the outward
contour thus has just $u=pZ_h,Z_h/c$. At D the extra scales stay on
the inside of this outward contour. The inward computation for a
higher positive interchanges H and D, giving $u=psZ_h,sZ_h/c$ at D;
its one-body zero is at $u=Z_h$. This exhausts the mixed term and
proves the lemma.
\end{proof}

\subsection{Summing the residues}
Record the step of each positive and each negative. An inverted pair
selects the same step for both signs. At an H step only this pair can
supply a positive, and at a D step only this pair can supply a negative.
For a choice with ordinary shifts, the assigned spectral values are
\[
  s p^t Z_h\quad\hbox{on the positive side},\qquad
  p^t \tau_h\quad\hbox{on the negative side},\qquad 0\le t<m
\]
at a step of multiplicity $m$ for the indicated sign. In the additional
choice of Lemma~\ref{lem:cubic-poles} replace $p$ by $1/c$ on the
single displaced value at its primary step.

\begin{lemma}\label{lem:cubic-pin}
In the signed sum over all cuts, a nonzero assignment must have as
many positives at $i$ as negatives at $i-1$ for every $i$. In this
case all the fields with these indices lie at their common center $i$.
\end{lemma}

\begin{proof}
We first identify the dependence on the cuts. At a primary step $h$
split $W_h(k,j)=U_k^h V_j^h$ into the products over $k\le t<h$ and
$h<t<j$, extending $U_{h+1}^h=V_h^h=0$. Give the positive the factor
$U_k^h$ and the negative $V_j^h\Phi_{h,j}$. The two transport ratios are
\[
 U_{i+1}^h/U_i^h=r_{Z_h}(Z_i),\qquad
 (V_{i-1}^h\Phi_{h,i-1})/(V_i^h\Phi_{h,i})=r_{Z_{h+1}}(Z_i).
\]
For the second ratio a horizontal primary gives forward factors
$R(v)$ or $R(v)/S(v)$ through H or D, with $v=Z_h/Z_{i-1}$.
A downward primary gives $R(v/c)$ or $R(v)$. Substitution of the
two possible steps proves the identity, including its cleared endpoint.
For all other one-body factors the same relations follow by dividing
$P_{i+1}/P_i$ or $N_{i-1}/N_i$ at the assigned spectral value.

The cut-independent factors are the one-body weights at the two outer
cuts, the primary residue constants and the mixed kernels at their
assigned values (with the primary kernel omitted). What remains is
\eqref{eq:cubic-density} with its transports, applied to the labels
$p^t Z_h$ and $p^t Z_{h+1}$ of either sign. To see that all allocations
occur with their stated coefficients, temporarily distinguish all
slots of each sign. A fixed primary at $h$ requires exactly the strict
inversion $k\le h<j$. Once its two slots are assigned, these endpoint
conditions are independent. For ordinary fields each allowed single
pair occurs once in the three inversions of the preceding proof. Two
disjoint inversions necessarily use $a,d$, giving the single double
difference. Their four mutual contractions have just the product of
the two same-sign and two mixed kernels, as in that proof. For an
inverted second power sum the two possible singleton choices and the
simultaneous elementary choice have precisely the coefficients of
$P(J)$ by \eqref{eq:cubic-mixed}. An allocation in a single bilinear
insertion cannot straddle a primary step, since its negative cut is
at most its positive cut. Thus including all three partitions of
three in the density adds no spurious primary terms. The factorials
of indistinguishable insertions are the $m_j!$ in
\eqref{eq:cubic-density} after summing the labeled allocations.
These arguments also cover the grouped term, which has no primary.
An ordinary unshifted cutoff is enforced by its zero transport at the
physical end of its interval.

For an ordinary repeated selection, specialize the labels of a length
$m$ string to $z,pz,\ldots,p^{m-1}z$. First divide the distinct-label
density by the ordered kernels of the consecutive copies, and take
the limit with independent consecutive ratios. There is one simple
pole on each such divisor for decreasing levels in positive cut order
and increasing levels in negative cut order. Reverse order and
membership in a single elementary subset are regular. Nonconsecutive
copies in a common subset cannot survive either, because the missing
intermediate copy must be strictly between them. All other same-sign
factors are regular, including the pair at distance two. Hence the
limit keeps exactly the ordered $p$-shifts; the difference between the
divided kernel and its logarithmic residue is a cut-independent
scalar. For a step carrying a primary, at most one sign is repeated.
Its upper mixed kernel is regular for an ordinary $p$-shift.
Equivalently this procedure reproduces successive singleton shifts of
one Macdonald variable, while an elementary subset shifts distinct
variables only. The same-sign orders are unaffected by the inserted
mixed kernels. The physical support of each upper slot is automatic,
since its unshifted positive is used last, or its unshifted negative first.

For the additional choice at H or D no same-sign limiting operation is
needed. The mixed kernel at that pole is replaced by its logarithmic
residue, $(1-1/p)(1-s)/(1-s/p)$ at H and its negative at D. The
opposite same-sign order has zero factor
$K_{1/p,1/s}(c)=0$ at H or $K_{p,s}(1/c)=0$ at D. Consequently
the base negative precedes the additional negative at H, and the
additional positive precedes the base positive at D, unless both
are in one simultaneous elementary subset. In the latter case the
two share their cut and retain $E(J)$. The label of the additional
slot is the label of its base divided by $c$. Transporting it throughout
the interval adds no choices: its base has the genuine endpoint zero
and the opposite strict order vanishes. The formula is therefore the
same distinct-label density, specialized at regular same-sign
arguments, with the same endpoint constraints.

Now extract the initial positive and final negative one-body weights
in a finite window containing all selected steps and their successor
vertices. In this window the cut sum is invariant under reordering
the parameters $Z_i$ by Lemma~\ref{lem:cubic-slices}. This remains
true after the repeated-label limits. Transport denominators cannot
vanish at a physical slice: if $Z_i$ and the base label have the
same $x$-degree, their ratio is one, before the specified powers of
$p$ or $1/c$ are inserted. Suppose more positives occur at $i$ than
negatives at $i-1$, and bring $Z_i$ to the first slice. Its unshifted
positive cannot pass this slice because $r_{Z_i}(Z_i)=0$. All its
upper positives must occur there as well. A negative based at another
vertex cannot be used here, since the backward transport of its base
would cross its own zero and the base must occur first. There are too
few available negatives for equal degrees of both signs on this slice.
If any multiplicity disagrees there is such an excess somewhere.
When all multiplicities agree, proceed in physical order. At the first
occupied vertex all its positives must have been used by this cut
and only its own negatives can yet be used. Their common center and
native cut are forced. Removing these slots and proceeding to the
next vertex proves the assertion.
\end{proof}

The ordinary matched residues have a simple interpretation. Consecutive
occupied contents are joined across their shared primary step. An
occupied interval opens at H and closes at D and consists of the
southeast-most box on each intervening content. A multiplicity greater
than one requires an H birth and a D death, hence is possible only at
a removable corner. Conversely, at size three each layer contains
at most one box on each diagonal: two on the same diagonal force a
square of four. The connected components of the bottom layer are
therefore border strips with successive contents, and an upper box
at a nonremovable corner would force its neighbor to have height
at least two, requiring at least four boxes in total. This gives the bijection between
the ordinary matched patterns and the reverse plane partitions of
size three.

\subsection{Weights of the surviving heights}
We compare their full weights. For a layer at content $i$ and height
$t$ counted from zero put $z=p^tZ_i$. Its positive and negative
variables have values $sz,z/x$, even when one of them was fixed
by fusion with a neighboring layer. For two such layers in increasing
content order their four regular cross contractions have product
\begin{equation}\label{eq:cubic-cross}
 \begin{split}
 \mathcal K(y)={}&K_{1/p,1/s}(y)K_{p,s}(y)
          K_{1/p,s}(sxy)K_{p,1/s}(y/(sx))\\
       ={}&R\big((p^{-1}-1)(1-x)(y+y^{-1})\big),\qquad
             y=z_{\rm later}/z_{\rm earlier}.
 \end{split}
\end{equation}
In the last expression $R$ acts multiplicatively on a character. The
identity follows by expanding the four binomial quotients. It gives
exactly the cross term of the quadratic height polarization between
the two layers, whose characters are $c/z_{\rm earlier}$ and
$c/z_{\rm later}$.

At a single corner of content $i$ a stack of $k\le3$ copies can use
only $e_1^j$ in each Newton power, since an elementary subset cannot
shift the same root twice. Its scalar is
\[
 h_k[l]=[v^k]\exp\!\sum_{j\ge1} l_j v^j/j
       =\prod_{t=1}^k(1-cp^{t-1})/(1-p^t).
\]
The last equality follows by comparing the generating function at $v$
and $pv$. The selected H and D variables give opposite finite Cauchy
shifts, which cancel their mutual bridge. Each layer $z=p^tZ_i$ has
the exterior factor
\begin{equation}\label{eq:cubic-outside}
 \prod_{b<i-1,\ {\rm H}}\frac{R(z/(xZ_b))}{R(z/Z_b)}
 \prod_{d>i,\ {\rm D}}\frac{R(z/Z_d)}{R(cz/(xZ_d))}.
\end{equation}
Indeed the horizontal Macdonald and downward Cauchy singleton
factors at an earlier $b$ give $s^{-1} R(z/(xZ_b))/R(z/Z_b)$,
and the two factors at a later $d$ give
$sR(z/Z_d)/R(cz/(xZ_d))$. The counts of the two types cancel
the negative normalization $s^i$ for each layer. On the one-box
diagram the height weighs $s^{-k}\prod_{t=1}^k R(p^{-t})=h_k[l]$.
On a general diagram each exterior step contributes to the linear
part of the polarization separately for each layer, giving
\eqref{eq:cubic-outside}. This proves the equality for a single stack;
\eqref{eq:cubic-cross} now treats disjoint stacks and interactions
between separate border strips.

For completeness we compute the one-layer weight of a connected
three-box strip. Write $\beta,\alpha$ for its opening H and closing
D, $b=Z_\beta$, $d=Z_\alpha$. Its internal H and D step weights are
denoted by $u$ and $v$, respectively. Delete the strip, leaving
$\mu\subset\lambda$. If $Q$ is the sum of its box characters and
$E_\mu$ the sum of those of $\mu$, reflection
$R(y)R(1/(sy))=s$ in the polarization gives its weight
\[
 R\big(\big[(1-x)(1-c/x)E_\mu-c\big]Q^*\big).
\]
In fact the unreflected increment is
$Q/p+(1-x)(E_\mu Q^*-QE_\mu^*/p)$; the scalar from
reflecting the first term cancels $s^{-|\,\lambda/\mu\,|}$.
Each step outside $[\beta,\alpha]$ gives
\eqref{eq:cubic-outside}, with $e<\beta$ for the H steps and $f>\alpha$ for the D steps in its
corresponding product, for each of the three layer parameters $z$.
On the residue side the opening and closing singletons furnish
$R(Z_\beta/Z_e)/R(Z_\alpha/Z_e)$ at an earlier H step $e$ and
$R(Z_\alpha/Z_f)/R(cZ_\beta/Z_f)$ at a later D step $f$.
Multiplication by the intervening $\Phi$ factors gives the same
exterior products by the two step recursions for $Z_i$.

It remains to compare on the isolated strip. The two endpoints and
their mutual kernel reduce to the scalar $l_1$, since
$l_1\gamma a_0=-1$. Every internal H leaves $1/R(cd/u)$ and every
internal D leaves $1/R(c^2 b/v)$ after cancellation against its
endpoint contractions. For an H preceding a D their two factors
from $\Phi$ and the mutual contraction $U_+(pt)U_-(t)$,
$t=v/(cu)$, have product $S(v/u)S(cu/v)$. These identities follow
on substituting the displayed $U_e$ and $S$. Thus the character
whose image under $R$ gives the residue is
\begin{equation}\label{eq:cubic-ribbon-patch}
 -c-\sum_u cd/u-\sum_v c^2b/v+
       (1-c)\sum_{u\ {\rm before}\ v}(v/u+cu/v).
\end{equation}
One can check the height expression on the four possible patches
directly. Put $Y=d+\sum u+\sum v$. Removing the strip divides the
intervening vertex weights by $c$ and exchanges the initial H and
final D. Telescoping $(1-x)(1-c/x)E_\mu-c$ along the horizontal
steps of the remaining patch gives the height character
$c\big((x^{-1}-1)\sum u^{-1}-1/d\big)Y$.
Substitution of the two step ratios gives the following common
values for this character and \eqref{eq:cubic-ribbon-patch}:
\[
\begin{array}{c|l}
 {\rm HH}&-c(1+x+x^2)\\
 {\rm HD}&-c-x^2-c^2/x^2+(1-c)(x+c/x)\\
 {\rm DH}&-c-cx-c^2/x\\
 {\rm DD}&-c(1+c/x+c^2/x^2).
\end{array}
\]
This also exhibits the mutual factor on the strip with internal
word HD. The analogous one-internal-letter calculation is the
domino of the preceding section.

There remains the domino with a doubled corner. At this corner
the first negative and the last positive are the base fields. Keep
the other two variables independent while taking the residue of the
base domino. Taking their same-sign residues next shifts each
spectrum by $p$, as in the single stack. The mixed kernel with
the neighbor stays regular: its later-to-earlier ratio in
$K_{p,1/s}$ is $p/s$ on HHD and $1/p^2$ on HDD. The local
factor at the corner is thus the ratio $W_2/W_1$ of the stack
and singleton weights already calculated. The upper layer still
interacts with the neighbor, giving \eqref{eq:cubic-cross}. Hence
the result in the two cases is
\[
 W_{\rm domino}\frac{W_2}{W_1}\,\mathcal K(y),\qquad
 y=px\quad({\rm HHD}),\qquad y=x/(cp)\quad({\rm HDD}).
\]
The quadratic height character gives exactly this expression by
adding the upper layer. Together these calculations account for
every ordinary matched pattern.

\begin{proof}[Completion of the proof of Theorem~\ref{thm:cubicdisk}]
For an additional pole of Lemma~\ref{lem:cubic-poles},
Lemma~\ref{lem:cubic-pin} leaves only a doubled removable end of a
domino. At that center one sign has an ordinary $p$-string. On the
other sign take $(q,V,t)=(1/p,1/s,1/c)$ on H and $(p,s,c)$ on D,
where $t$ is the ratio of the later to the earlier value in the
permitted order. For two ordinary insertions the unique order on
the $p$-string and the zero of the reverse-order kernel on the other
sign leave just $l_1^2 K_{q,V}(t)/2$. For one second power sum the
ordinary string again allows just its singleton order, while the
other sign contributes its two-variable Newton factor. The sum,
after removal of their common factors, is
\[
 \tfrac12\big[(l_1^2+l_2)K_{q,V}(t)-2l_2 E_V(t)\big]=0 .
\]
Indeed in both cases $K_{q,V}(t)/E_V(t)=
2l_2/(l_1^2+l_2)$ by substitution. This removes all the additional
residues. The ordinary matched residues give the complete height sum,
which proves the theorem.
\end{proof}

For example on the two-box row the contribution of the additional
pole in the mixed summand, including its coefficient in
\eqref{eq:cubic-moment}, is
\[
 \frac{c(c+1)p(x+1)(x-1)^2(c-1)^2(x-c^2)}
 {2(1-p)(x-pc)(px-c)(pc-1)(px^2-c)} .
\]
At $c=3,x=5$ its residue and constant coefficient in $1-p$ are
$-576/11$ and $-16704/121$. Both are canceled by the ordinary cube.
